% Source: Wald GR, physical PDF pp. 19--21
%         (printed pp. 6--8).
% Coverage: Section 1.3, The Spacetime Metric; equation (1.3.1).
% Figures/tables/footnotes: no new figures, tables, or footnotes; source float
%                           Figure 1.3 is transcribed with Section 1.2.
% Status: source-reviewed and compile-clean.
% Uncertainties: none.

\subsection{The Spacetime Metric}\label{sec:1.3}

In the previous section, we gave a prescription for how an inertial observer \(O\) can
label the events in spacetime with global inertial coordinates \(t, x, y, z\). However,
a fundamental tenet of special relativity is that there are no preferred inertial
observers. As seen above, a different inertial observer using the same procedure
assigns different labels \(t', x', y', z'\) to the events in spacetime. Thus, the
coordinate labels themselves do not have intrinsic significance since they depend as
much on which observer does the labeling as they do on the properties of spacetime
itself. It is of great interest to determine what quantities have absolute,
observer-independent significance, i.e., truly measure intrinsic structure of
spacetime. This is equivalent to determining what functions of global inertial
coordinates are independent of the choice of inertial frame.

In prerelativity physics the answer is the following. The time interval \(\Delta t\)
between two events has absolute significance; all observers will agree on the value of
\(\Delta t\). Furthermore, the spatial interval \(\lvert\Delta\mathbf{x}\rvert\) between two
simultaneous events is observer independent. However, these quantities (or functions
of them) are the only ones with absolute significance. For example, observers moving
with nonzero relative velocity will disagree over the spatial interval between
nonsimultaneous events.

In special relativity neither the time interval nor the space interval between
``relatively simultaneous'' events (i.e., events determined to be simultaneous by a
particular observer) has absolute significance. The quantity which is observer
independent is the spacetime interval, \(I\), defined by

\begin{equation}
  I = -(\Delta t)^2
      +\frac{1}{c^2}\bigl[(\Delta x)^2+(\Delta y)^2+(\Delta z)^2\bigr].
  \tag{1.3.1}\label{eq:1.3.1}
\end{equation}

Indeed, the Poincar\'e transformations (the set of all possible transformations
between global inertial coordinates) consist precisely of the linear transformations
which leave \(I\) unchanged. The spacetime interval \(I\) and functions of \(I\) are
the only observer-independent quantities characterizing the spacetime relationships
between events.

What is truly remarkable about the expression for \(I\) is that it is quadratic in the
coordinate differences, just like the distance function in Euclidean (i.e., flat,
positive definite) geometry. Indeed, the only difference is the minus sign in front of
\((\Delta t)^2\), allowing \(I\) to become zero or negative. We shall refer to \(I\) as
the metric of spacetime in analogy to an ordinary Euclidean metric. (More precisely,
the metric of spacetime in special relativity will be defined later to be a tensor field
associated with the formula for the spacetime interval between two ``infinitesimally
nearby'' events; see Eq.~\eqref{eq:4.2.2} below.) As we shall see in chapter 3, this
difference in metric signature makes very little difference in the mathematical
analysis of metrics. In particular definitions of geodesics (``straightest possible
lines'') and curvature carry through in the same way for metrics with the signature
of \(I\) as for ordinary positive definite metrics. It is interesting to note that, as
discussed more fully in chapter 4, the paths in spacetime of inertial observers in
special relativity are geodesics of the spacetime metric, and the curvature associated
with \(I\) is zero, i.e., the spacetime geometry in special relativity is flat.
